\subsection{强互素多项式}

设 R 为交换环。若 R 的两个元素 x、y 的主理想 Rx、Ry 互素，则称 x、y 强互素，换言之（第二章，第 1 节，第 2 号），若 \(\mathbf{R}x + \mathbf{R}y = \mathbf{R}\)；这等价于存在两个元素 \(a, b\)（属于 \(\mathbf{R}\)），使得 \(ax + by = 1\)。

引理 1（“Euclid 引理”）。设 \(x, y\) 是 \(\mathbf{R}\) 中的两个强互素元素；若 \(z \in \mathbf{R}\) 满足 \(x\) 整除 \(yz\)，则 \(x\) 整除 \(z\)。

若 \(1 = ax + by\)，则 \(z = x(az) + (yz)b\)。

若 \(x\) 和 \(y\) 在 \(\mathbf{R}\) 中强互素，则

\[
\mathrm{R} x y = (\mathrm{R} x) \cap (\mathrm{R} y)
\]

（第二章，第 1 节，第 2 号，命题 5）；若 R 是整环，则两个强互素元素的最小公倍元等于它们的乘积（代数，第六章，第 1 节，第 8 号），因而按照代数第六章第 1 节第 12 号的意义它们互素。反之，若 R 是主理想整环，则两个互素元素也强互素，这由 Bézout 等式得出（代数，第七章，第 1 节，第 2 号，定理 1）。

对于多项式环，有如下结果：

命题 1。设 A 为交换环，P、P' 是 A{[}X{]} 中的两个强互素多项式。假设 P 是首一的，次数为 s。则 A{[}X{]} 中的每个多项式 T 都可以唯一写成如下形式：

\[
\mathrm{T} = \mathrm{PQ} + \mathrm{P} ^ {\prime} \mathrm{Q} ^ {\prime}\tag{1}
\]

其中 \(\mathbf{Q} \in \mathbf{A}[\mathbf{X}]\)、\(\mathbf{Q}' \in \mathbf{A}[\mathbf{X}]\)，且 \(\deg(\mathbf{Q}') < s\)。

若进一步有 \(\deg (\mathbf{T})\leqslant t\) 且 \(\deg (\mathbf{P}^{\prime})\leqslant t - s\)，则 \(\deg (\mathbf{Q})\leqslant t - s\)

由于 \(\mathbf{P}\) 是首一的，\(\mathrm{PR} \neq 0\) 对每个多项式 \(\mathbf{R} \neq 0\)（属于 \(\mathbf{A}[\mathbf{X}]\)）都成立，且此时 \(\deg(\mathrm{PR}) = s + \deg(\mathbf{R})\)。

设 \(\mathbf{T}\) 是 A{[}X{]} 中任意多项式。由于 P 和 \(\mathbf{P}'\) 生成的理想就是整个 A{[}X{]}，存在多项式 \(\mathbf{Q}_1\) 和 \(\mathbf{Q}_1'\)，使得

\[
\mathrm{T} = \mathrm{PQ} _ {1} + \mathrm{P} _ {1} ^ {\prime} \mathrm{Q} _ {1} ^ {\prime};
\]

由于 P 是次数为 s 的首一多项式，Euclid 除法（代数，第四章，第 1 节，第 5 号）表明存在两个多项式 Q'、Q'\,'，使得 Q'\_1 = PQ'\,' + Q'，其中 deg(Q') \textless{} s；于是得到

\[
\mathrm{T} = \mathrm{PQ} _ {1} + \mathrm{P} ^ {\prime} (\mathrm{PQ} ^ {\prime \prime} + \mathrm{Q} ^ {\prime}) = \mathrm{PQ} + \mathrm{P} ^ {\prime} \mathrm{Q} ^ {\prime}
\]

其中 \(Q = Q_{1} + P'Q''\)。为证明公式（1）的唯一性，只需证明关系式

\[
0 = \mathrm{PQ} + \mathrm{P} ^ {\prime} \mathrm{Q} ^ {\prime}, \quad \deg (\mathrm{Q} ^ {\prime}) <   s\tag{2}
\]

蕴含 \(Q = Q' = 0\)。现在，若（2）成立，则 \(P\) 整除 \(-PQ = P'Q'\)，又由于 \(P\) 和

\(P'\) 强互素，根据引理 1，\(P\) 整除 \(Q'\)；若 \(Q' \neq 0\)，则存在多项式 \(S \neq 0\)，使得 \(Q' = PS\)，从而

\[
\deg (\mathbf {Q} ^ {\prime}) = s + \deg (\mathbf {S}) \geqslant s,
\]

这就矛盾了。我们得到 \(Q' = 0\)，于是 PQ = 0，最后根据开头的注记得到 Q = 0。

最后，假设 \(\deg(T) \leqslant t\) 且 \(\deg(P') \leqslant t - s\)；将多项式 \(T\) 写成形式（1）后，

\[
\deg (\mathrm{P} ^ {\prime} \mathrm{Q} ^ {\prime}) \leqslant \deg (\mathrm{P} ^ {\prime}) + \deg (\mathrm{Q} ^ {\prime}) <   s + \deg (\mathrm{P} ^ {\prime}) \leqslant t
\]

因而

\[
s + \deg (Q) = \deg (P Q) = \deg (T - P ^ {\prime} Q ^ {\prime}) \leqslant t
\]

所以 \(\deg (\mathbb{Q})\leqslant t - s.\)

例。对于多项式 \(P \in A[X]\)，P 与 X - a 强互素（其中 \(a \in A\)）的充要条件是 \(P(a)\) 在 A 中可逆。因为若 P 与 X - a 强互素，则由命题 1，存在 \(c \in A\) 及多项式 \(Q \in A[X]\)，使得 \(cP + (X - a)Q = 1\)，从而 \(cP(a) = 1\)，且 \(P(a)\) 可逆。反之，由 Euclid 除法

\[
\mathbf {P} = (\mathbf {X} - a) \mathbf {R} + \mathbf {P} (a)
\]

并且，若 \(\mathrm{P}(a)=b^{-1}\)，其中 \(b\in A\)，则得到 \(1=b\mathrm{P}-b(\mathrm{X}-a)\mathrm{R}\)，这表明 P 与 X-a 强互素。

设 A、B 为两个交换环，\(f\colon\mathbf{A}\to\mathbf{B}\) 为环同态。若 \(\mathbf{P}=\sum_{i\geqslant0}a_{i}\mathbf{X}^{i}\) 是 \(\mathbf{A}[[\mathbf{X}]]\) 中的形式幂级数，令 \(\bar{f}(\mathbf{P})\) 表示形式幂级数 \(\sum_{i\geqslant0}f(a_{i})\mathbf{X}^{i}\)，它属于 \(\mathbf{B}[[\mathbf{X}]]\)。若 \(\mathbf{P}\) 是多项式，则 \(\bar{f}(\mathbf{P})\) 也是多项式；若进一步 \(\mathbf{P}\) 是首一的，则 \(\bar{f}(\mathbf{P})\) 也是首一的，且次数与 \(\mathbf{P}\) 相同。最后，\(\mathbf{P}\mapsto\bar{f}(\mathbf{P})\) 显然是从 \(\mathbf{A}[[\mathbf{X}]]\) 到 \(\mathbf{B}[[\mathbf{X}]]\) 的同态，它延拓 \(f\)，并将 \(\mathbf{X}\) 映到 \(\mathbf{X}\)。在本段其余部分中，我们将始终以此意义使用记号 \(\bar{f}\)。

命题 2。设 A、B 为两个交换环，f 为从 A 到 B 的同态，P、P' 为 A{[}X{]} 中的两个多项式。若 P 与 P' 在 A{[}X{]} 中强互素，则 \(\bar{f}(\mathbf{P})\) 与 \(\bar{f}(\mathbf{P}')\) 在 B{[}X{]} 中强互素。反之，若 \(f\) 是满射，且其核包含于 A 的 Jacobson 根中，并且 P 是首一的，则结论也成立。

假设 P 与 \(P'\) 强互素；则存在多项式 Q、\(Q'\)，属于 A{[}X{]}，使得 \(PQ + P'Q' = 1\)；从而得到

\[
\bar {f} (\mathbf {P}) \bar {f} (\mathbf {Q}) + \bar {f} (\mathbf {P} ^ {\prime}) \bar {f} (\mathbf {Q} ^ {\prime}) = 1,
\]

由此得到第一个断言。为证明第二个断言，令 a 表示 f 的核；

令 \(\mathbf{E} = \mathbf{A}[\mathbf{X}]\)，令 \(\mathbf{F}\) 为 \(\mathbf{E}\) 中由 \(\mathbf{P}\) 和 \(\mathbf{P}'\) 生成的理想；由于 \(f\) 是满射且 \(\bar{f}(\mathbf{P})\) 首一，命题 1 表明，对于每个多项式 \(\mathbf{T} \in \mathbf{A}[\mathbf{X}]\)，存在两个多项式 \(\mathbf{Q}, \mathbf{Q}'\)，属于 \(\mathbf{A}[\mathbf{X}]\)，使得

\[
\bar {f} (\mathrm{T}) = \bar {f} (\mathrm{P}) \bar {f} (\mathrm{Q}) + \bar {f} (\mathrm{P} ^ {\prime}) \bar {f} (\mathrm{Q} ^ {\prime}),
\]

于是有关系 \(E = F + aE\)。现在，E/F 是有限生成的 A-模，因为每个多项式模 P 同余于一个次数 \(<\deg(P)\) 的多项式，而 P 是首一的。由于 \(E/F = a(E/F)\)，且 a 包含于 A 的 Jacobson 根中，Nakayama 引理表明 \(E/F = 0\)（代数，第八章，第 6 节，第 3 号，命题 6 的推论 2），这意味着 P 与 \(P'\) 强互素。

